\chapter{Background}
\label{ch:background}

\section{The GPU memory hierarchy}

The arithmetic units of a modern GPU are far faster than the path that feeds
them. The reference card sustains \TiledTflopsAtMaxN{} TFLOP/s of single precision
arithmetic on this workload, by the convention Chapter~\ref{ch:results} states, while its
device memory delivers \GpuMeasuredBandwidth{} GB/s in practice. At four bytes per float
that is about \MeasuredFloatsBillions{} billion floats per second, so a
kernel that performs one operation per float loaded runs at under one
percent of the arithmetic the hardware has shown it can do. Every technique in this
project is a way of raising the number of operations performed per byte
brought in from device memory.

Three levels of the hierarchy matter here.

\paragraph{Coalescing.} A warp of thirty two threads issues its loads together.
When those threads read consecutive addresses, the hardware merges them into a
small number of wide transactions. When they read scattered addresses it issues
many narrow ones and discards most of what it fetched. This is why the particle
state is stored as a structure of arrays rather than an array of structures: if
one particle's seven floats sat together, a warp reading the $x$ component of
thirty two consecutive particles would touch thirty two separate cache lines
and use a fraction of each. With separate arrays the same read is thirty two
consecutive floats \cite{cudaguide}.

\paragraph{Shared memory.} Each multiprocessor has a small block of memory,
\SI{\GpuSharedMemoryKib}{\kibi\byte} per block on this card, that all threads of a block can
read at a latency and bandwidth far better than device memory. A block can
cooperatively stage a tile of data there once and then have all its threads
read it many times. This is the mechanism behind the tiled force kernel.

\paragraph{Registers and the load store path.} Values held in registers cost
nothing to read. Between registers and device memory sit the first level cache
and the load store pipeline, and as Chapter~\ref{ch:results} shows, on this
problem that pipeline rather than device memory is what actually limits the
naive kernel.

\section{Gravitational N-body dynamics}

The acceleration of body $i$ under the gravity of all the others is

\begin{equation}
  \mathbf{a}_i = G \sum_{j \neq i} m_j
  \frac{\mathbf{r}_j - \mathbf{r}_i}
       {\left(\left|\mathbf{r}_j - \mathbf{r}_i\right|^2 + \varepsilon^2\right)^{3/2}}.
  \label{eq:gravity}
\end{equation}

The softening length $\varepsilon$ in the denominator is not a numerical
convenience to be apologised for. Without it the acceleration diverges as two
bodies approach, a single time step throws the pair apart at enormous speed,
and the run is destroyed. For a collisionless system, where the smooth mean
field rather than individual close encounters carries the dynamics, softening
is the physically appropriate model \cite{aarseth2003}. The shipped runs are not fully in
that regime, as Chapter~\ref{ch:discussion} notes, and there softening is what keeps the
integrator stable. The potential must be softened consistently,

\begin{equation}
  U = -G \sum_{i < j} \frac{m_i m_j}
      {\sqrt{\left|\mathbf{r}_j - \mathbf{r}_i\right|^2 + \varepsilon^2}},
  \label{eq:potential}
\end{equation}

because pairing a softened force with an unsoftened potential breaks energy
conservation for a reason that has nothing to do with the integrator. Both come
from the same expression in the code, and a test differentiates
Equation~\ref{eq:potential} numerically and compares it against
Equation~\ref{eq:gravity}.

\subsection{Why velocity Verlet}

The integrator advances

\begin{align}
  \mathbf{r}_{n+1} &= \mathbf{r}_n + \mathbf{v}_n \Delta t
                      + \tfrac{1}{2}\mathbf{a}_n \Delta t^2, \\
  \mathbf{v}_{n+1} &= \mathbf{v}_n
                      + \tfrac{1}{2}\left(\mathbf{a}_n + \mathbf{a}_{n+1}\right)\Delta t.
\end{align}

Three alternatives were considered and rejected.

Explicit Euler is first order and not symplectic. Its energy error grows
without bound, and on a circular orbit it spirals outward.

Fourth order Runge Kutta is accurate per step but requires four force
evaluations per step. Since the force evaluation is essentially the entire cost
of this simulation, that is a factor of four in exchange for accuracy that a
symplectic method achieves differently.

Velocity Verlet is second order, needs exactly one force evaluation per step,
and is symplectic: it exactly conserves a modified Hamiltonian near the true
one, so its energy error oscillates within a bound rather than accumulating
\cite{verlet1967,hairer2006}. The single precision energy figure in
Section~\ref{sec:gates} is too dominated by rounding to show this, so the order is tested
in double precision.

The fused kernel uses the leapfrog form,

\begin{align}
  \mathbf{v}_{n+1/2} &= \mathbf{v}_{n-1/2} + \mathbf{a}_n \Delta t, \\
  \mathbf{r}_{n+1}   &= \mathbf{r}_n + \mathbf{v}_{n+1/2} \Delta t,
\end{align}

which produces the same position sequence to rounding, provided the velocity is
offset backwards by half a step once before the loop begins. That offset, and
the fact that the stored velocity afterwards is not synchronised with the
positions, is a detail with real consequences for snapshots and for the energy
diagnostic; both are handled explicitly in Chapter~\ref{ch:implementation}.

\subsection{The Plummer sphere}

Benchmarks need a physically sensible initial condition. A uniform random cloud
is not one: it is far from equilibrium and collapses immediately, so it
measures the wrong thing. The Plummer sphere \cite{plummer1911} has density

\begin{equation}
  \rho(r) = \frac{3M}{4\pi a^3}\left(1 + \frac{r^2}{a^2}\right)^{-5/2},
  \qquad
  \frac{M(r)}{M} = \frac{r^3}{\left(r^2 + a^2\right)^{3/2}},
\end{equation}

and the cumulative profile inverts in closed form, so radii are sampled
directly. Speeds are drawn from the exact distribution function by the
rejection method of Aarseth, H\'{e}non and Wielen \cite{aarseth1974}, which puts
the sphere in virial equilibrium at $t = 0$. Drawing speeds from a convenient
Maxwellian instead would leave it out of equilibrium, and it would visibly
breathe on the first crossing time in a way that looks exactly like an
integrator bug.

\section{Smoothed particle hydrodynamics}

SPH represents a continuous field by interpolating over particles with a
smoothing kernel $W$ of finite support $h$ \cite{monaghan1992}. Density is

\begin{equation}
  \rho_i = \sum_j m_j W\!\left(\left|\mathbf{r}_i - \mathbf{r}_j\right|, h\right),
\end{equation}

and pressure follows from a stiff equation of state $p = k(\rho - \rho_0)$,
which keeps the fluid weakly compressible without solving a pressure Poisson
equation every step.

Three different kernels are used, each for the one job it does well
\cite{muller2003}. The \emph{poly6} kernel is a polynomial in $r^2$, so density
needs no square root. The \emph{spiky} kernel supplies the pressure gradient,
because its gradient does not vanish as the separation goes to zero; the poly6
gradient does, which means no restoring force exactly where the fluid needs one
most, and particles pile up on top of each other. The \emph{viscosity} kernel
has a positive Laplacian throughout its support, so the viscous term always
damps and never injects energy. Each integrates to one over its support, which
is the property the normalisation gate checks by quadrature.

\subsection{Finding neighbours}

The neighbour set changes every step and must be rebuilt every step. Every fluid
domain here is bounded and has walls, so the search uses a dense uniform grid, with cells at
least as wide as the smoothing radius, and a neighbour can then only be in one of the
twenty seven cells around a particle. Cell $(i,j,k)$ of an $n_x \times n_y \times n_z$ grid has
the index

\begin{equation}
  c(i,j,k) = i + n_x \left(j + n_y\,k\right),
\end{equation}

and particles are sorted by it. Two things follow. The index is unique, so two visited cells
can never be the same slot, and no neighbour can be counted twice. And $x$ varies fastest, so
the three cells of an $x$ row are contiguous in sorted order, and the twenty seven cells are nine
ranges. Section~\ref{sec:sphresults} reports the measured agreement with a brute force sum,
which is what proves it.

Earlier versions used spatial hashing \cite{teschner2003} with a table of twice as many slots as
particles, on the argument that a dense grid would waste its storage on empty cells. That
argument was wrong, as Chapter~\ref{ch:corrections} records: on these domains the table was
larger than the grid it avoided. A k-d tree is a poor fit for a full rebuild on a GPU every step.
